\subsubsection{One finite group for all primes different from the characteristic}

\begin{lemma}[Prime-uniform generation by arithmetic unit cups]
\label{lem:B-unit-cup-prime-uniform}
\status{Proved conditional on divisor Tate for surfaces}
Let \(k\) be a finite field of characteristic \(p\), and assume
divisor Tate for all smooth projective surfaces over finite
extensions of \(k\), at one fixed prime \(\ell_0\ne p\).
Put \(D=(\Q/\Z)[p']\).  For every smooth quasi-projective
\(k\)-scheme \(T\), arithmetic cyclic cups give a map
\[
 c_T:\Gamma(T,\mathcal O_T^\times)\otimes D
       \longrightarrow \Br(T)[p']
\]
with finite cokernel.  The cokernel is one finite group across
all primes different from \(p\); separate finiteness of its
primary components is not substituted for this assertion.
\end{lemma}
\begin{proof}
First suppose that \(U\) is a dense open of a smooth projective
geometrically connected \(Y/k\), and put
\(\Gamma=\operatorname{Gal}(\bar k/k)\).  Let \(B\) be the
free group on the geometric prime divisors in \(Y\setminus U\).
Regularity gives the exact sequence
\[
 0\longrightarrow\bar k^\times
 \longrightarrow E=\Gamma(\bar U,\mathcal O_{\bar U}^\times)
 \longrightarrow B\longrightarrow\operatorname{Pic}(\bar Y)
 \longrightarrow\operatorname{Pic}(\bar U)\longrightarrow0.
\]
Boundary subsets of codimension at least two do not change these
units or Picard groups.  The lattice
\(N=E/\bar k^\times\subset B\) is finite free, with finite
\(\Gamma\)-action.

All positive cohomology of \(\bar k^\times\) under
\(\Gamma\) vanishes: in degree one this is Hilbert~90, or
surjectivity of \(x\mapsto x^{\#k-1}\), and in higher degrees
it follows from torsion cohomological dimension one.  Hence
global units of \(U\) surject onto \(N^\Gamma\), and
\[
 H^2(\Gamma,E)\simeq H^2(\Gamma,N)
 \simeq H^1(\Gamma,N\otimes\Q/\Z)
 \simeq N_\Gamma\otimes\Q/\Z.
\]
For the middle isomorphism, higher continuous cohomology of the
discrete rational representation \(N\otimes\Q\) vanishes by
averaging over finite quotients.  The last isomorphism is the
procyclic torsion-cohomology formula.  Since \(N\) has finite
Galois action, \(N^\Gamma\to N_\Gamma\) has finite cokernel.
Tensoring with the divisible group \(\Q/\Z\) makes it
surjective.  The same holds on the prime-to-\(p\) summand.
The induced map is precisely the arithmetic cyclic cup: its
degree-two cocycle is the Bockstein of the arithmetic character
multiplied by the invariant unit.  Consequently the image of
\(H^2(\Gamma,E)[p']\) in \(H^2(U,\Gm)\) is contained
in \(C_U=\operatorname{im}(c_U)\).

The other term needed in total degree two is finite over all
primes as well.  Let \(L\) be the image of \(B\) in
\(\operatorname{Pic}(\bar Y)\), and let
\(P=\operatorname{Pic}^0_{Y/k,\mathrm{red}}\).
The group \(P(\bar k)\) is torsion, so
\(F=L\cap P(\bar k)\) is finite.  With
\(Q=\operatorname{NS}(\bar Y)/\operatorname{im}(L)\),
there is an exact sequence
\[
 0\longrightarrow P(\bar k)/F
 \longrightarrow\operatorname{Pic}(\bar U)
 \longrightarrow Q\longrightarrow0.
\]
Frobenius minus identity on \(P\) is an isogeny, since its
differential is minus identity.  It is therefore surjective
on geometric points and \(H^1(\Gamma,P(\bar k))=0\).
Together with \(H^2(\Gamma,F)=0\), this gives
\(H^1(\Gamma,P(\bar k)/F)=0\).  The group \(Q\) is
finitely generated with finite Galois action, and
\(H^1(\Gamma,Q)\) is finite: its finite torsion part is
harmless, and the free part reduces to finite-group cohomology,
because a profinite group has no nonzero continuous homomorphism
to a discrete torsion-free lattice.  Thus
\(H^1(\Gamma,\operatorname{Pic}(\bar U))\) is finite.

The Hochschild--Serre filtration on the kernel of
\[
 H^2(U,\Gm)\longrightarrow H^2(\bar U,\Gm)^\Gamma
\]
has a quotient of \(H^2(\Gamma,E)\) and a subquotient of
\(H^1(\Gamma,\operatorname{Pic}(\bar U))\) as its two
associated groups.  Every \(H^2(\Gm)\) group here is torsion
because the schemes are regular.  The preceding calculations
therefore show that the prime-to-\(p\) part of this kernel,
modulo \(C_U\), is finite.  Arithmetic cups vanish geometrically.
No degeneration or splitting of the spectral sequence is used.

The remaining geometric input is prime-uniform.  Surface Tate,
the Kahn--Morrow reduction, and the unconditional divisor
splitting of Lemma~\ref{lem:B-divisor-splitting} imply
\(\Br(\bar U)^\Gamma[\ell_0^\infty]\) is finite by
geometric Kummer and finite cotype.  Qin's Theorem~5.2 then
implies
\[
                 \Br(\bar U)^\Gamma[p']\quad\text{is finite}.
\]
This theorem applies to smooth geometrically connected opens;
it does not require an additional full Frobenius-semisimplicity
conjecture~\cite[Theorem~5.2 and its proof in
Lemmas~3.2--3.4, Propositions~4.2--4.3, and Theorem~4.4]
{BQin2021Uniform}.
For precision concerning versions, the splitting and
prime-independence results cited in that proof as an earlier
Qin manuscript are Proposition~2.6 and Theorem~2.7 in the
opened August~2025 Ertl--Keller--Qin version
\cite[Proposition~2.6 and Theorem~2.7]{BErtlKellerQin2025}.
The arithmetic kernel calculation now proves that
\(\Br(U)[p']/C_U\) is finite.

For general geometrically connected smooth quasi-projective
\(T\), choose a projective generically finite alteration
\(g:U\to T\) of degree \(d\), where \(U\) is open in
a smooth projective scheme.  Apply de~Jong's alteration theorem
to a projective compactification; source and target have ample
line bundles~\cite[Theorem~4.1 and Remark~4.2]{c_dejong1996}.
For every \(n\) prime to \(p\), the finite-coefficient Gysin
map gives
\[
 g_*:H^2(U,\mu_n)\longrightarrow H^2(T,\mu_n),
 \qquad g_*g^*=d.
\]
The scalar identity does not require \(d\) invertible modulo
\(n\): the composite is an endomorphism of the constant sheaf,
and its value on the dense finite-flat locus is multiplication
by \(d\).  Riou's finite-flat normalization includes inseparable
generic degree~\cite[Proposition~2.5.13, Remark~2.5.14,
Corollary~3.1.2, and Lemmas~3.5.8--3.5.10]{c_riou2014}.

These maps descend to compatible maps on \(\Br[n]\).
Indeed, a Kummer class from \(\operatorname{Pic}(U)/n\)
restricts to zero at the generic field.  Its trace has zero
generic Brauer class, hence zero Brauer class on regular \(T\),
by injectivity into the function-field Brauer group
\cite[Proposition~1.4 and Corollary~1.8]{BGrothendieckBrauerII1966}.
Thus the induced trace is independent of the Kummer lift; no
unstated integral derived trace on \(\Gm\) is needed.
The projection formula and the Kummer norm identity give
\[
 \operatorname{cor}_g(\chi|_U,u)
             =(\chi,\operatorname{Nm}_{k(U)/k(T)}u).
\]
The norm of a global unit on \(U\) is a global unit on \(T\).
Properness gives centers for every extension of every
codimension-one valuation of normal \(T\); all orders of the
unit are zero, so their norm valuation is zero.  The degree-one
Kummer trace identity can be checked on the dense finite-flat
locus and extended by the injectivity of
\(H^1(T,\mu_n)\to H^1(T^\circ,\mu_n)\) on normal schemes.
This argument includes purely inseparable factors.

Pullback and trace therefore act on
\(B_T=\Br(T)[p']/C_T\) and \(B_U=\Br(U)[p']/C_U\),
with composite \(d\).  The first case, applied over the finite
constant field of \(U\), makes \(B_U\) finite; the character
rescaling in the final paragraph identifies its cup image with
the one used here.  Thus one
integer \(h>0\) annihilates \(B_T\).  The subgroup \(C_T\)
is divisible, so \(B_T[h]\) is a quotient of
\(\Br(T)[h]\), after removing the \(p\)-power factor of
\(h\).  Kummer and finite-coefficient cohomology make the latter
finite.  Hence \(B_T=B_T[h]\) is finite.

Finally, if a connected component has finite constant field
\(k'/k\), the character normalized over \(k\) restricts to
\([k':k]\) times the one normalized over \(k'\).
Multiplication by this integer on \(D\) is surjective, so
the cup images coincide.  Finite disjoint unions complete the
proof.
\end{proof}

\begin{theorem}[Whole prime-to-characteristic Brauer finiteness and the exact gate]
\label{thm:B-Brauer-prime-uniform}
\status{Proved equivalence; the surface-Tate condition remains unresolved}
Fix \(p\), put \(R=\Z[1/p]\) and \(D=(\Q/\Z)[p']\),
and for every finite-type \(\F_p\)-scheme \(S\) define
\[
 A_S=H^2_{\et}(S,\Gm)\otimes R,
 \qquad C_S=\operatorname{im}
 \bigl(\Gamma(S,\mathcal O_S^\times)\otimes D\to A_S\bigr).
\]
The following conditions are equivalent:
\begin{enumerate}[label=(\roman*)]
\item Divisor Tate holds, at one prime different from \(p\),
      for all smooth projective surfaces over all finite
      extensions of \(\F_p\).
\item \(A_S/C_S\) is a finitely generated \(R\)-module
      for every finite-type \(\F_p\)-scheme \(S\).
\item \(\Br(S)[p']/C_S\) is a finite group for every such
      \(S\).
\item The whole prime-to-\(p\) Brauer cokernel is finite for
      every proper smooth geometrically connected curve family
      over every such \(S\).
\end{enumerate}
Under these conditions the relative Brauer conclusion also holds
for every finite-type morphism identifying global units,
in particular every smooth proper morphism with geometrically
connected fibres, in arbitrary relative dimension.
All statements include singular, nonreduced, nonproper, and
nonseparated bases.
\end{theorem}
\begin{proof}
Assume (i).  Repeat the proof of
Theorem~\ref{thm:B-all-bases-Brauer-units}, now over
\(R=\Z[1/p]\).  We give the points that make this a single
prime-uniform argument.  Choose the same finite-type
\(h\)-cover \(q:T\to S\), with \(T\) a finite disjoint
union of smooth quasi-projective schemes.  The kernel
\[
 K=\ker(A_S\to A_T)
\]
is finitely generated over \(\Z[1/p]\) itself by
Lemma~\ref{lem:B-Gm-proper-kernel}; its proof used the
localized Picard and unit groups, not a prime-dependent bound.
Since \(T\) is regular, \(A_T=\Br(T)[p']\).
Lemma~\ref{lem:B-unit-cup-prime-uniform} supplies one integer
\(N\) annihilating \(A_T/C_T\).

The unit-cup kernel in Lemma~\ref{lem:B-unit-cup-kernel}
is finite for \(D\) as well: the proof uses one full-rank
integer valuation matrix.  Its Smith invariant factors have
finite prime-to-\(p\) torsion, supported at finitely many
primes.  Likewise the rationally exact sequence of finite unit
lattices for \(S,T,T_1\) has finite middle homology after
tensoring with \(D\).  Choose fixed integers \(M,L\)
annihilating these two finite groups.  The exact same cup and
\v Cech calculation gives, for every \(a\in A_S\), a
\(v\in\Gamma(S,\mathcal O_S^\times)\otimes D\) with
\[
                      LMN\,a-c_S(v)\in K.
\]
Thus the image of \(A_S/C_S\to A_T/q^*C_S\) is killed
by the fixed integer \(h=LMN\), and its kernel is a quotient
of \(K\).

One does not need, and does not assume, a finite-cotype
description summed over all primes.  Instead, \(q^*C_S\)
is divisible.  If \(hb\in q^*C_S\), choose
\(c\in q^*C_S\) with \(hc=hb\); then \(b-c\in A_T[h]\).
Consequently
\[
              A_T[h]\longrightarrow(A_T/q^*C_S)[h]
              \quad\text{is surjective}.
\]
The source is finite by Kummer and finite-coefficient cohomology,
using the prime-to-\(p\) part of \(h\).
The image of \(A_S/C_S\) is therefore finite, and its kernel
is finitely generated, proving (ii).

Localization cannot turn a nontorsion element into a torsion
element.  Hence \((A_S)_{\rm tors}=\Br(S)[p']\).
Since \(C_S\) is torsion, the torsion submodule of \(A_S/C_S\)
is exactly \(\Br(S)[p']/C_S\).  A finitely generated torsion
\(\Z[1/p]\)-module is finite, giving (iii).
If pullback identifies units, functoriality gives
\(f^*C_S=C_X\); its relative Brauer cokernel is a quotient
of \(\Br(X)[p']/C_X\), hence finite.  For the smooth proper
families in (iv), equality of units follows from the
base-change-compatible identity \(f_*\mathcal O_X=\mathcal O_S\)
\cite[Lemma~36.32.6, tag 0E0L]{BStacksConnectedUnits}.
This proves (iv) and the stronger concluding assertion.

Finally (iv) implies the fixed-prime relative-curve assertion,
which implies (i) by the pencil argument of
Theorem~\ref{thm:B-relative-surface-gate}.
Alternatively (iii) applied to a smooth projective surface has
\(C_S=0\), because its units are finite, and the proper
Kummer criterion again gives (i).  This completes the
equivalences.  The characteristic-primary counterexample in
Theorem~\ref{thm:B-Brauer-characteristic} is unaffected.
\end{proof}

\begin{corollary}[Prime-uniform abelian torsors and an integral lattice]
\label{cor:B-torsors-prime-uniform}
\status{Proved equivalence with the same unresolved surface-Tate condition}
Under the equivalent conditions of
Theorem~\ref{thm:B-Brauer-prime-uniform}, every abelian scheme
\(A/S\) over a finite-type \(\F_p\)-scheme satisfies
\[
 H^1_{\et}(S,A)[p']\ \text{finite},
 \qquad
 H^1_{\et}(S,A)\otimes\Z[1/p]
       \ \text{finitely generated over }\Z[1/p].
\]
Conversely, the universal first assertion implies the surface-Tate
condition.  These are statements about one whole prime-to-\(p\)
group, including all primes simultaneously.  They include
nonreduced and nonseparated bases and abelian schemes without
a chosen global polarization.
\end{corollary}
\begin{proof}
Let \(f:B=A^\vee\to S\), let \(e\) be its zero section,
and write \(P=\operatorname{Pic}_{B/S}=R^1f_*\Gm\) on the
\(\et\) site.  Relative Picard duality identifies
\(P^0\) with \(A\).  Here \(P^0\) means fibrewise algebraic
equivalence to zero on every test scheme; it remains the dual
abelian scheme over a nonreduced base.  In particular we use the
arbitrary-abelian-scheme form, not just the projective form of
the theorem~\cite[section~2.3.2, Remark~2.3.6 and Theorem~2.3.7]
{BRobert2023AbelianSchemes}.
Set \(N=P/A\) as an \(\et\)-sheaf; no local constancy of
\(N\) is asserted or needed.  Inversion on \(B\) acts as
\(-1\) on \(A\), by duality, and as \(+1\) on \(N\).
For the latter assertion, \([-1]^*L\otimes L^{-1}\) is
algebraically trivial on every geometric fibre; this is exactly
the defining fibrewise criterion for \(P^0\).  The underlying
line-bundle identities are the theorem-of-the-square identities
in Lemma~7.16(ii) and Corollaries~7.23(ii) and~7.25 of the
opened duality chapter~\cite{BMoonenDuality2011}.

Naturality of the connecting homomorphism in
\[
 0\longrightarrow A\longrightarrow P\longrightarrow N
       \longrightarrow0
\]
therefore gives
\[
 -\delta(n)=\delta([-1]^*n)=\delta(n),
 \qquad 2\delta(n)=0\quad(n\in H^0(S,N)).
\]
Consequently the kernel of \(H^1(S,A)\to H^1(S,P)\)
is killed by two.

The zero section splits the unit map in the derived category:
\[
       Rf_*\Gm\simeq\Gm\oplus K,
       \qquad K\in D^{\ge1}(S_{\et}),
       \qquad\mathcal H^1(K)=P.
\]
The truncation triangle
\(P[-1]\to K\to\tau_{\ge2}K\) has
\(\mathbb H^1(S,\tau_{\ge2}K)=0\).  Hence it induces
an actual injection
\[
 H^1(S,P)\lhook\joinrel\longrightarrow
 H^2(S,K)\lhook\joinrel\longrightarrow H^2(B,\Gm).
\]
Its image is in the kernel of \(e^*\), and therefore meets
\(f^*H^2(S,\Gm)\) trivially.  In particular torsion inputs
give actual torsion Brauer classes; no lift from a torsion
quotient of a nontorsion group is being assumed.  We obtain
\[
 H^1(S,A)[p']\longrightarrow
 \operatorname{coker}\bigl(\Br(S)[p']\to\Br(B)[p']\bigr)
\]
with kernel killed by two.  If \(p=2\), that kernel is zero.
If \(p\ne2\), it is finite because Kummer gives a surjection
\(H^1(S,A[2])\to H^1(S,A)[2]\), and the finite \(\et\)
sheaf \(A[2]\) has finite cohomology over the finite-type
finite-field scheme \(S\).  A finite affine open cover proves
the same finiteness when \(S\) is nonseparated.
The target is finite by
Theorem~\ref{thm:B-Brauer-prime-uniform}, since \(B\to S\)
is smooth proper with geometrically connected fibres.  This
proves the first assertion.

For the second assertion, repeat the proof of
Proposition~\ref{prop:B-integral-torsor-lattice} over
\(\Z[1/p]\).  Its actual \v Cech-cocycle map already
embeds the proper-cover restriction kernel, after inverting
\(p\), in a subquotient of a finitely generated
\(\Z[1/p]\)-module of abelian-scheme sections.  On the regular
alteration the torsor group is torsion; its localization is
now finite by the first assertion, simultaneously at every
prime different from \(p\).  The same extension argument,
bounded \(p\)-primary nilpotent comparison, and finite open-cover
descent therefore prove the stated finite generation.
Finally, the universal first assertion implies its fixed-prime
version, whose surface-Tate equivalence is
Theorem~\ref{thm:B-all-bases-surface-gate}.
\end{proof}
